\chapter*{Conclusiones}
\addcontentsline{toc}{chapter}{Conclusiones}

Este trabajo se propuso comparar dos enfoques diferentes para resolver una tarea clínica sensible como es la evaluación de la calidad de imágenes retinianas: un modelo entrenado específicamente (ResNet-18 fine-tuned) y un modelo multimodal de propósito general (GPT-4o) utilizado sin entrenamiento adicional, bajo modalidades zero-shot y few-shot.

Los resultados muestran que, en términos cuantitativos, el modelo entrenado ofreció un desempeño significativamente superior, tanto en precisión como en estabilidad. GPT-4o, en cambio, demostró una sensibilidad alta pero con una gran cantidad de falsos positivos, lo cual compromete su confiabilidad como herramienta autónoma en un contexto clínico.

Desde el punto de vista práctico, los modelos multimodales presentan ventajas importantes, especialmente en modalidad zero-shot, donde no requieren recolección ni etiquetado de datos ni conocimientos técnicos avanzados para su uso. En few-shot, si bien se requiere una selección previa de ejemplos etiquetados, el volumen de datos necesarios sigue siendo considerablemente menor en comparación con un proceso completo de entrenamiento supervisado. Sin embargo, este menor esfuerzo técnico inicial no elimina la necesidad de conocimiento experto: al momento de diseñar el prompt o elegir las imágenes de ejemplo, se hace evidente que el desempeño del modelo depende en gran medida de la claridad y relevancia clínica de estas decisiones. Aún así, los modelos multimodales siguen representando herramientas de bajo esfuerzo y rápida implementación. No obstante, este mismo grado de flexibilidad puede volverse difícil de controlar: el diseño del prompt y la selección de ejemplos demostraron tener un impacto directo en la predicción, lo cual reintroduce un sesgo que los modelos de redes neuronales intentan minimizar: el sesgo del investigador. En este sentido, el prompt engineering funciona casi como una nueva forma de "ingeniería de features".

También debe considerarse el aspecto económico. Si bien el costo de uso de GPT-4o fue accesible en esta tesis (aproximadamente 17 USD para procesar 10.000 imágenes en modalidad few-shot), sigue siendo un modelo ofrecido como servicio. Cada clasificación implica una llamada a la API, cuyo costo varía según el tamaño del input y la longitud de la respuesta. En contraste, los modelos como ResNet-18 pueden descargarse gratuitamente y ejecutarse localmente, aunque el fine-tuning requiere recursos de cómputo cuyo costo, aunque presente, es más predecible y bajo control del desarrollador.

Además, el uso de modelos ofrecidos como servicio, como GPT-4o, plantea desafíos para la reproducibilidad y la trazabilidad: el usuario no tiene control sobre los pesos del modelo, ni garantía de estabilidad entre versiones. Esto dificulta su integración en sistemas clínicos donde la transparencia y la validación son críticas. También deben considerarse aspectos relacionados con la seguridad, como las alucinaciones (respuestas inventadas pero plausibles), y la falta de entrenamiento específico en patologías concretas como la retinopatía diabética, lo cual limita su capacidad para ajustar su criterio al contexto clínico real.

En conclusión, si bien los modelos multimodales como GPT-4o muestran un potencial interesante como herramientas auxiliares de bajo esfuerzo, su aplicación actual en tareas clínicas automatizadas requiere precaución. Son sistemas versátiles y configurables, pero aún carecen de la robustez, precisión y trazabilidad necesarias para ser confiables en escenarios sensibles como el de la salud. Su uso futuro dependerá tanto de avances técnicos como de la capacidad para integrar estas herramientas en flujos de trabajo controlados, validados y auditables.